% ════════════════════════════════════════
% 摘要（00.摘要.tex）
% ════════════════════════════════════════
\begin{abstract}

小区微网的购电计划必须在负载与光伏尚未发生时制定，这段"先计划、后兑现"的时间差是本题的核心约束。我们以购电费用最小为目标，用两阶段随机规划与滚动时域优化方法逐一建模求解。

\textbf{问题一}中电价与负载固定、光伏预测已知，属单日确定性优化。把一天按 10 分钟步长离散为 144 个时段，以购电功率、充放电功率和储电量为决策变量，建立功率平衡、储能递推、容量与功率上下限约束，并令 0:00 与 24:00 储电量同为 6000 kWh。该线性规划含 577 个连续变量，可精确求解，得全天购电 \textbf{59482.70 kWh}、购电费 \textbf{35126.95 元}；低谷充电、高峰放电是费用下降的主要来源。

\textbf{问题二}的难点在于每天 0:00 负载与光伏都尚未发生，直接引用附件二的实际值就是用事后信息做决策。我们以 1 月为训练样本，用扩张窗口逐日标定预测模型，在 0:00 给出当天 144 个间隔的功率预测，再由历史残差估计不确定性，负载与光伏各取低、中、高三档组合成 9 个情景，建立两阶段\textbf{随机线性规划}，目标为计划购电费与 5 倍电价下期望紧急购电费之和。运行时储能严格按计划充放电，真实功率到达后按缺口结算紧急购电、富余电量计为弃电。以 1 月为预热期，按附件 5 口径统计 2 月 1 日至 12 月 31 日共 334 天：计划购电 \textbf{21833065.92 kWh}、计划购电费 \textbf{13325494.59 元}，紧急购电 \textbf{282910.68 kWh}、紧急购电费 \textbf{1057398.40 元}，总购电费 \textbf{14382892.99 元}，日均约 43062.55 元；334 天中有 333 天出现紧急购电，紧急电量约占计划电量的 1.30\%。

\textbf{问题三}可在 0:00、6:00、12:00、18:00 获得未来 24 小时整点光伏预报。把小时级预报展成 10 分钟序列，0:00 制定全天计划购电，其余预报时刻对未执行时段滚动重优化；计划高于调整的部分按 0.5 倍电价计违约费，调整高于计划的部分按 1.5 倍计价，分段费用用上境变量线性化。含调整的总费用为 \textbf{14337999.15 元}，比只用 0:00 计划降低 \textbf{444073.36 元}，降幅约 \textbf{3.00\%}。

\textbf{问题四}中电价实时波动，我们在附件四电价下重算问题二与问题三：计划购电费 \textbf{13895479.01 元}、紧急购电费 \textbf{1181652.65 元}，合计 \textbf{15077131.66 元}；问题三总费用 \textbf{14956526.18 元}，比只用 0:00 计划节省 \textbf{120605.48 元}。波动电价整体抬高购电成本，滚动调整可对冲其中一部分。

本文有两点创新。其一，以"先预测、再决策、后结算"的三段式框架组织四问，计划购电只依赖 0:00 可得的信息，杜绝用事后数据做决策；其二，用历史残差构造情景并建立两阶段随机规划，以 5 倍紧急电价刻画缺电风险，再由滚动时域优化吸收逐次到达的预报。四问最终都化为连续变量线性规划，决策变量全部连续，全局最优性可核验，且模型仅依赖电价与不可控出力两列输入，替换数据即可复用。

\keywords{\textbf{微网}；\textbf{储能调度}；\textbf{随机规划}；\textbf{滚动时域优化}；\textbf{光伏预报}}

\end{abstract}
